\section{Related Work}
The landscape of Network Intrusion Detection Systems (NIDS) has evolved significantly in recent years, propelled by the availability of comprehensive datasets such as CICIDS2017. A substantial body of literature over the past 3-5 years has explored various data-driven approaches; however, effectively handling extreme multi-class imbalance remains a critical challenge. This section reviews recent advancements, categorized into traditional Machine Learning (ML), Deep Learning (DL), and Ensemble learning approaches.

\subsection{Traditional Machine Learning Approaches}
Traditional ML algorithms, including Support Vector Machines (SVM), Decision Trees (DT), Naive Bayes, and k-Nearest Neighbors (k-NN), have been extensively applied to NIDS. Recent studies demonstrate their mathematical simplicity and low computational overhead during inference [1], [2]. However, these solitary classifiers typically exhibit a high sensitivity to specific feature spaces and scale poorly with the high-dimensional, high-velocity nature of modern network traffic. More critically, traditional ML algorithms inherently prioritize the majority class to maximize global accuracy [3]. Consequently, they systematically misclassify minority attack classes (e.g., Infiltration or Heartbleed in CICIDS2017) as benign traffic, leading to an unacceptable rate of False Positives and missed critical intrusions.

\subsection{Deep Learning Approaches}
To capture complex, non-linear relationships within network traffic, researchers have increasingly adopted Deep Learning architectures. Convolutional Neural Networks (CNNs) and Deep Autoencoders have shown superior capability in automated feature extraction without requiring extensive manual engineering [4], [5]. Furthermore, recurrent architectures such as Long Short-Term Memory (LSTM) networks have been leveraged to analyze temporal traffic patterns [6]. Despite their robust representational power, DL models are notoriously data-hungry and prone to severe overfitting when trained on the highly skewed distributions typical of modern network datasets [7]. The minority classes are often treated as noise during the gradient descent optimization process. Additionally, training these deep networks entails substantial computational cost and hyperparameter sensitivity, making them difficult to optimize and deploy efficiently in real-time environments.

\subsection{Ensemble Methods}
Recognizing the limitations of individual classifiers, ensemble learning has emerged as a state-of-the-art paradigm in cybersecurity. Homogeneous ensemble techniques, such as Random Forest, use bagging to reduce model variance but are constrained by the shared inductive biases of their base decision trees [8]. Heterogeneous and boosting-based ensembles, incorporating models like XGBoost and LightGBM, have proven highly effective due to their sequential error-correction mechanisms and computational efficiency on tabular data [9], [10]. However, existing ensemble literature in NIDS frequently overlooks the compounding effect of extreme data imbalance. Most current studies either apply algorithms directly to raw datasets or rely on simplistic, noisy data-augmentation techniques (e.g., basic SMOTE) that distort the high-dimensional feature space, failing to adequately isolate and detect rare attack vectors [11], [12].

\subsection{Addressing the Research Gap}
The fundamental limitation across the reviewed literature is the failure to comprehensively construct an architecture that natively and synergistically addresses both algorithmic variance and data-level extreme imbalance. Current studies often treat sampling and classification as disparate steps, resulting in sub-optimal detection of minority attacks.

Our proposed framework explicitly fills this research gap. By engineering a novel \textit{Hybrid Ensemble Architecture}—which strategically utilizes the structural precision of XGBoost and LightGBM as base learners and the advanced non-linear mapping capabilities of an MLP meta-classifier—we break the dependency on a single algorithmic paradigm. Crucially, this algorithmic diversity is tightly coupled with our advanced \textit{Balanced Sampling} strategy directly tailored for the CICIDS2017 dataset. This integrated pipeline ensures that the feature space is mathematically rebalanced, amplifying the signal of minority attacks without injecting stochastic noise. Consequently, our approach comprehensively mitigates the algorithmic bias prevalent in solitary models, drastically reducing False Positives (FPs) while maintaining exceptionally high recall for stealthy intrusions.
